\titlespacing*{\subsection}{0pt}{1.0em}{0.3em}
\titlespacing*{\paragraph}{0pt}{1.0em}{0.3em}

본 연구의 핵심 목표는 환율 예측 결과를 단순히 수치 예측에 그치지 않고, 멀티 에이전트 시스템을 통해 경제적 해석과 인사이트를 제공하는 AI 분석 프레임워크를 구축하는 데 있다. 이를 위해 멀티 에이전트 시스템은 단순한 텍스트 도출을 넘어 전문화된 에이전트 간 순차적 추론 구조를 채택하였다. 이 구조는 실제 금융기관의 전문가 역할을 모사하여 5단계의 전문가 에이전트 파이프라인으로 구성된다. 먼저, 계량경제 에이전트가 선행 예측 모델의 산출물인 정량적 수치와 불확실성을 통계적으로 해석하면, 이를 바탕으로 거시경제 에이전트와 통화정책 에이전트가 순차적으로 개입하여 수치에 숨겨진 환율의 변동 요인을 설명한다. 이와 같이 통화정책 에이전트는 금리와 같은 통화 기조가 환율 변화에 미치는 영향을 설명한다. 이후 비정형 위험 에이전트가 거시·정책적 시나리오에 기반한 잠재적 비정형 요소의 리스크를 검증하고, 최종적으로 종합 에이전트가 앞선 모든 논의를 종합하여 시장 인사이트를 도출한다. 이와 같은 선형적 형태의 구조는 에이전트가 단편적 관점에 매몰되는 것을 방지한다. 또한 개별 에이전트는 RAG의 지식 베이스를 활용하여 앞선 에이전트의 분석을 비판적으로 수용하고 확장함으로써, 블랙박스로 여겨지던 예측 모델의 산출물에 명확한 경제학적 원인과 논리적 근거를 설명하는 중요한 역할을 수행한다.

\subsection{환율 상호작용의 구조적 특성과 네트워크 기반 예측 해석}
실험 결과에서 관측된 가장 주목할 만한 특징은 세 통화가 동일한 방향으로 움직이지 않는다는 점이다. 미국 달러 인덱스(TWDI)가 비교적 강한 흐름을 유지하는 동안, 한국 원화와 일본 엔화는 서로 다른 속도와 불확실성 범위를 가지며 약세 흐름을 형성하였다. 이는 국제 금융 환경이 비슷하더라도 각국의 통화 정책 여건과 거시경제 상황에 따라 통화 간 상호 변동 경로가 서로 다르게 나타날 수 있음을 보여 준다. 이러한 차이는 기존의 단일 시계열 기반 모델이나 두 통화 중심의 분석에서는 충분히 반영되기 어려우며, 실제로 국가 간 상호작용을 제한적으로 설명하는 경향이 있다. 
반면, 본 연구에서 채택한 B--MTGNN은 다변량 시계열 전체를 하나의 그래프 구조로 표현하여 한 국가의 변화가 다른 통화에 어떤 방식으로 영향을 미치는지 함께 학습한다. 국가 간 환율 격차($\mathrm{Gap}_{i,j}$) 분석에서 일본--한국 조합의 격차가 상대적으로 크게 나타난 결과는 통화 간 변동이 단순한 상관관계만으로 설명되지 않음을 보여 준다. B--MTGNN의 그래프 학습 모듈은 사전에 정의된 관계 없이도 데이터로부터 통화 간 잠재적 의존 관계를 스스로 학습한다. 이렇게 학습된 그래프 구조는 국가 간 경제적 연결이나 정책 환경의 차이를 반영하며, 예측 경로의 형성에 직접 관여한다. 특히 일본의 신뢰구간이 상대적으로 넓게 나타난 결과는 일본 통화의 변동성이 높음을 시사하며, 이는 B--MTGNN이 국가 간 상호 영향을 함께 고려하여 예측 경로를, 형성했음을 보여 준다.

\subsection{예측 및 위험 분석에 대한 시사점: 멀티 에이전트 해석}
환율 예측 모델은 과거 패턴을 연장하여 미래를 예측하므로, 예상치 못한 대외적 변수에 의한 시장 변화를 포착하기 어렵다. 이러한 한계를 보완하기 위해 다중 에이전트 시스템을 도입하였다. 계량경제 분석가가 비선형적 변동을 이벤트 신호로 식별하면, 거시경제 분석가와 통화정책 전략가가 이를 금리차 및 정책 기조의 변화 관점에서 해석하고, 비정형 위험 분석가가 지정학적 변동과의 연결 고리를 규명한다. 이러한 구조는 의사결정자에게 예측 근거에 대한 맥락 정보를 제공하여 실용적인 설명 가능한 예측을 구현한다. 특히 2026년 예측치와 점진적으로 확대되는 신뢰구간을 단순한 통계적 오차로 표현하는 것이 아니라, 각 분야 전문 에이전트들의 교차 검증을 통해 통화정책과 외생적 리스크가 복합적으로 작용한 결과임을 성공적으로 입증하였다. 결론적으로, 본 시스템은 단순한 텍스트 생성이 아닌 금융 시장의 복합적 요소를 심층적으로 식별하고 해석할 수 있는 타당성을 지니며 이를 통해 단일 모델 대비 환각과 확증 편향을 감소시켜 전문성을 고도화하였음을 시사한다\cite{FinVision2024}.

\subsection{방법론적 기여 및 일반화 가능성}
본 연구의 방법론적 기여는 크게 세 가지로 요약된다. (1) 적응형 그래프 학습을 통해 사전에 정의되지 않은 통화 간 상호의존성을 데이터로부터 직접 추출하였고, (2) 몬테카를로 드롭아웃(MC Dropout) 기반의 베이지안 추론을 통해 단순 점 추정을 넘어선 확률적 예측으로 확장하였다. (3) 정량적 모델의 출력에 LangGraph와 LightRAG 기반의 다중 에이전트 시스템을 결합하여, 수치적 결과에 경제적 맥락을 부여하는 차세대 AI 분석 프레임워크를 구현하였다.
분석 대상이 미국, 한국, 일본 3개국 네트워크에 한정된 것은 데이터의 단순한 한계가 아니라, 외환 시장의 변동 구조의 본질적 메커니즘을 정밀하게 표현하기 위한 의도적 설계이다. 무분별한 노드 확장은 오히려 모델의 학습 성능을 저하할 위험이 있다. 따라서 정책 투명성이 높고 신뢰할 수 있는 공식 지표(Fed, BOK, BOJ 등)가 보장된 핵심 3개국을 선정하여, 모델의 수치적 불확실성이 에이전트의 거시경제적 추론으로 정확히 입증되는지 그 타당성을 엄밀히 검증하였다. 추후에는 RAG 지식 베이스의 데이터와 지식 갱신 주기를 체계화하고, LLM 기반 멀티 에이전트 시스템으로 확장하며, 유로화(EUR) 등 대규모 다국적 경제권으로 본 프레임워크의 전세계적으로 적용 가능성을 입증하는 것이 후속 연구 과제로 남는다.
그럼에도, 본 연구가 제안하는 그래프 기반 확률적 예측과 멀티 에이전트 해석의 통합 구조는 환율 분야에 국한되지 않고, 비대칭적 상호작용이 존재하는 여타 금융 자산이나 다변량 시계열 예측 문제로 확장 적용이 가능하다. 결론적으로, 본 연구의 의의는 경제적 맥락 해석을 갖춘 차세대 AI 분석 프레임워크의 구조적 설계 원칙을 제시했다는 점에 있다.
